% TOP_PROBLEM: {"id": 3700042, "title": "Modified Cartan conjecture", "category_id": 8, "category": {"id": 8, "name": "analysis", "display_name": "Analysis", "description": "Limits, continuity, calculus, and function theory.", "slug": "analysis", "order_index": 8, "created_at": "2026-07-31T15:26:25.671Z"}, "status": "open"}
% FIRST_POSTED: null
% ENTRY_KIND: statement_only
% Imported from: batch12.tex; UnsolvedMath ID 3700042
% Compile twice: lualatex 3700042.tex
\documentclass[11pt]{article}
\usepackage{fontspec}
\setmainfont{Latin Modern Roman}
\usepackage{amsmath,amssymb,mathtools,unicode-math}
\setmathfont{Latin Modern Math}
\usepackage{xurl,hyperref}
\hypersetup{hidelinks,pdftitle={UnsolvedMath Problem 3700042}}
\newcommand{\sref}[2]{\hyperlink{source:#1:#2}{[S#2]}}
\newcommand{\cataloguescope}[1]{\paragraph{Mathematical question.}#1}
\begin{document}
% UnsolvedMath ID 3700042; source catalogue code AMR-036-0042
\setcounter{section}{3700041}
\section{Modified Cartan conjecture}\label{3700042}
{\small\sffamily UnsolvedMath ID 3700042 · Analysis}
\subsection{Definitions and mathematical statement}

\paragraph{Shared notation.}$\mathbb N=\{1,2,\ldots\}$, and $\mathbb Z,\mathbb Q,\mathbb R,\mathbb C$ denote the integers, rationals, reals and complex numbers. Any other conventions are specified locally.


Write $D(r)=\{z\in\mathbb C:|z|<r\}$ and $D=D(1)$. For $p\ge3$, consider a sequence of vectors of holomorphic zero-free functions
\[
 (f_{1,n},\ldots,f_{p,n}),\qquad \sum_{j=1}^{p}f_{j,n}\equiv0\quad(n\in\mathbb N).
\]
For a region $\Omega\subset D$, a nonempty set $I\subseteq\{1,\ldots,p\}$ is a $C$-class if there is a fixed index $k\in I$ such that all quotients $f_{j,n}/f_{k,n}$, $j\in I$, are locally uniformly bounded on $\Omega$, and
\[
 \sum_{j\in I}\frac{f_{j,n}}{f_{k,n}}\longrightarrow0
 \quad\text{locally uniformly on }\Omega\text{ as }n\to\infty.
\]
The dominant index $k$ does not vary with $n$. Each class has at least two indices. For a subsequence, apply this definition to the restricted sequence.

Let $R_p^{\mathrm{opt}}$ be the supremum of the radii $r\in(0,1]$ for which every such sequence has a subsequence admitting a partition of all $p$ indices into $C$-classes on $D(r)$.
\paragraph{Statement recorded in the supplied source.}
Prove $R_p^{\mathrm{opt}}>0$ for every $p\ge3$. Can one radius $R>0$ work for every $p$ (equivalently, is $\inf_{p\ge3}R_p^{\mathrm{opt}}>0$)? Give a geometric interpretation of the partition property.

For this interpretation, the vectors define holomorphic curves in
\[
 X_p=\{[x_1:\cdots:x_p]\in\mathbb P^{p-1}:\ \sum_jx_j=0,\ x_j\ne0\ \forall j\},
\]
the complement of $p$ hyperplanes in $\mathbb P^{p-2}$. For a block $I$, its coordinate projection to $\mathbb P^{|I|-1}$ should converge after extraction to a holomorphic curve in the hyperplane $\sum_{j\in I}x_j=0$; the source's Proposition 6.2 states the equivalence with the $C$-class partition.

Eremenko, Xu and Zhang's September 2026 Theorem 1.2 states $R_p^{\mathrm{opt}}>0$, and Theorem 1.3 gives $R_5^{\mathrm{opt}}=2-\sqrt3$. Those assertions alone do not settle the question of a positive radius uniform in $p$. \sref{3700042}{2}
\subsection{Short English statement}
Partition a subsequence of every zero-sum vector sequence into groups that asymptotically cancel on a disk of radius depending only on the number of functions. Also ask whether one positive radius works in all dimensions and interpret the grouping as convergence of projective coordinate projections.
\subsection{Sources}
\begin{description}
\item[\hypertarget{source:3700042:1}{S1}] ULAM.AI and the UnsolvedMath Contributors, UnsolvedMath: A Curated Collection of Open Mathematics Problems, record 3700042 (AMR-036-0042). CC BY 4.0. \url{https://huggingface.co/datasets/ulamai/UnsolvedMath}.
\item[\hypertarget{source:3700042:2}{S2}] Alexandre Eremenko, Zongben Xu and Teng Zhang, The modified Cartan conjecture (2026), Definition 1.1; Theorems 1.2–1.3; Proposition 6.2. \url{https://arxiv.org/abs/2609.23674}.
\end{description}
\label{end:3700042}
\end{document}
